\scanpage{61}
% archive sheet number: 60
quitte à passer à une extension finie de $k_0$, $\sigma\leadsto({}^{\sigma}g)^{-1}g$ \deleted{\unsure{ou}} \inserted{inverse de $g^{-1}g^\sigma$}, où $g\in G(\deleted{\unsure{\mathbb C}}k)$ est gén. /$k_0$ \inserted{(point déterminé mod translation par $g_0\in G(k_0)$)}. On trouve alors pour $g^\sigma$ \uline{tt} pt gén de $G(k)$, et par suite $g^\sigma g^{-1}$ peut \deleted{tt} \inserted{\unsure{exhauster|exister|exact}} les valeurs $a$ telles que $\deleted{\unsure{g}}ag$ soit gén /$k_0$, par ex. ttes valeurs telles que que \deleted{\unsure{$(ag)$|$(\sigma)$}} \deleted{\unsure{$ag$}} $a$ soit gén. sur $k_0(g)$.
\deleted{En d'autres \unsure{termes}} Considérons enfin l'hom total $\sigma\longrightarrow\sigma_*$
\diagram{\begin{tikzcd}[ampersand replacement=\&]
\operatorname{Gal}(k/k_0)\arrow[r,"{\deleted{\unsure{\varphi_{\infty}}}\inserted{\unsure{\varphi_v|\phi_v}}}"]\&G(\mathbb A_{\mathbb C})
\end{tikzcd}}

\deleted{\unsure{disons} On \unsure{a envie} de dire que l'image} Considérons le noyau de $\operatorname{Gal}(k/k_0)\unsure{\mathrel{\text{---}}|\longrightarrow}G(\mathbb A_{\unsure{\mathrm{fin}|\mathrm f}})$, qui contient le noyau $\operatorname{Gal}(k/\bar k_0)$ de $\operatorname{Gal}(k/k_0)\unsure{\longrightarrow|\mathrel{\text{---}}}\operatorname{Gal}(k/\bar k_0)$. Ce noyau est envoyé dans $G(\mathbb A_{\mathbb C})$ par $\varphi_{v\infty}$, et l'image est encore dense.
\deleted{De façon plus précise,} \inserted{Il en résulte que} \deleted{\unsure{$\varphi_v$ est}} \deleted{dominant} $\overline{\varphi_v(\operatorname{Gal}(k/k_0)}$ est un sous \unsure{(-groupe)|(-gpe)|(groupe)} ouvert de $G(\mathbb A_{\mathbb C})$.
\margin{(\phantom{o})?}
